\documentclass[11pt]{article}
% autocompile publish
\usepackage{math-hse}
\title{Лекция 14. Элементы теории вероятностей, часть 3}

\date{03.12.2025}
\begin{document}
\section{Формула полной вероятности}
Рассмотрим следующую ситуацию. Допустим, в некотором вузе есть три факультета: математический, физический и химический. Факультеты отличаются по числу студентов — например, на математическом факультете учится 1/10 студентов всего вуза, на физическом — 3/10, а на химическом — оставшиеся 6/10 (то есть 3/5) всех студентов вуза (см. табл.~\ref{otlichniki}).\footnote{Конечно, мы считаем, что каждый студент учится только на одном факультете, и каждый учится на каком-нибудь факультете из этих трёх.} Мы также знаем, какова доля отличников на каждом факультете. Например, на математическом факультете каждый четвёртый является отличником, на физическом — каждый второй студент отличник, а на химическом $3/4$ студентов являются отличниками.  Спрашивается, какова доля отличников в целом по вузу?

\begin{table}[hbt]
	\begin{center}
		\begin{tabular}{|c|c|c|c|c|c|c|}
			\hline
			факультет & матфак & физфак & химфак \\
			\hline
			какая часть студентов учится на факультете & 1/10  & 3/10 & 6/10 \\
			\hline
			доля отличников на факультете & 1/4 & 1/2 & 3/4\\
			\hline
		\end{tabular}
		\caption{Распределение студентов по факультетам}\label{otlichniki}
	\end{center}
\end{table}

\begin{figure}[hbt]
	\begin{center}
		\includegraphics[scale=0.7]{fp.pdf}
		\caption{Распределение студентов по факультетам. Закрашенные прямоугольники соответствуют отличникам. Их высоты — условным вероятностям.}\label{fig:fp}
	\end{center}
\end{figure}
Допустим для простоты, что всего в вузе учится 1000 студентов. Тогда на математическом факультете будет $1000\cdot \frac{1}{10}=100$ студентов. Среди них четвертая часть (то есть $1000\cdot \frac{1}{10} \cdot \frac{1}{4}=100\cdot \frac{1}{4}=25$ человек) являются отличниками, и это составляет $25/1000=0.025$ от студентов всего вуза. Нетрудно видеть, что это число не зависит от того, сколько студентов всего учится в вузе — если их $N$ человек, то среди них $N\cdot \frac{1}{10}$ учится на математическом факультете, отличников с математического факультета $N\cdot\frac{1}{10}\cdot\frac {1}{4}$ человек, и их доля среди всех $N$ студентов вуза равна $N\cdot\frac{1}{10}\cdot\frac{1}{4} \cdot \frac{1}{N}=\frac{1}{10}\cdot\frac{1}{4}=0.025$. (Вспомните задачи про фонды и компании, которые владели какими-то частями друг друга…)

Здесь также можно привести простую геометрическую интерпретацию (см. рис.~\ref{fig:fp}): на рисунке, отличникам с математического факультета соответствует крайний левый закрашенный прямоугольник, ширина которого равна $1/10$, а высота~--- $1/4$; общая доля отличников с математического факультета таким образом равна площади этого прямоугольника, то есть $0.025$.

Иными словами, можно записать:
\begin{multline}\label{eq:doli}
	\text{Доля отличников с мат. факультета среди студентов всего вуза} = \\
	= \text{доля студентов этого факультета среди студентов всего вуза } \times \\
	\times \text{ доля отличников на математическом факультете}
\end{multline}

На физическом факультете учится $1000\cdot \frac{3}{10}=300$ студентов. Отличников среди них половина, то есть $300\cdot \frac{1}{2}=150$ человек. Это составляет $150/1000=0.15$ от всех студентов вуза (средний закрашенный прямоугольник на рисунке). Наконец, на третьем — химическом — факультете учится $1000\cdot \frac{6}{10}=600$ студентов. Из них три четверти отличники, то есть $600\cdot \frac{3}{4}=450$ студентов, что составляет $450/1000=0.45$ от всех студентов вуза (крайний правый закрашенный прямоугольник).

Поскольку каждый отличник учится ровно на одном факультете, чтобы найти число всех отличников, достаточно сложить полученные числа: всего отличников $25+150+450=625$ человек. Доля отличников по всему вузу составляет $625/1000=0.625$.  Заметим, что к тому же результату мы бы пришли, если бы складывали \emph{доли отличников от всего вуза} по каждому факультету: $0.025+0.15+0.45=0.625$. Такую долю составляет закрашенная часть большого прямоугольника на рисунке~\ref{fig:fp} от всего большого прямоугольника. Ответ не изменился бы, если бы в вузе училось не 1000 студентов, а любое другое число (проверьте!).

Переведем теперь рассмотренную задачу на язык теории вероятностей. Допустим, мы выбираем случайного студента вуза. Пространство элементарных исходов здесь совпадает со множеством всех студентов (по-прежнему будем считать, что их 1000 человек). Можно рассмотреть различные события — например, событие $H_1$ — «студент учится на математическом факультете», событие $H_2$~--- «студент учится на физическом факультете», событие $H_3$~--- «студент учится на химическом факультете», а также событие $A$~--- «студент является отличником». Каковы вероятности таких событий?

Событию $H_1$ («студент учится на математическом факультете») благоприятствует столько исходов, сколько студентов учится на матфаке. Вероятность события $H_1$ в таком случае равна отношению числа студентов математического факультета к числу студентов всего вуза, то есть будет в точности совпадать с долей студентов мат. факультета во всём вузе — по условию, она равна 1/10. Аналогично обстоит ситуация с событиями $H_2$ и $H_3$.

Событию $A$ («студент оказался отличником») благоприятствует столько исходов, сколько отличников учится в вузе (в примере — 625 человек). Таким образом, вероятность случайно выбрать отличника из всего вуза — это в точности доля студентов-отличников во всем вузе (то есть 0.625).

Допустим теперь, что мы выбрали случайного студента среди студентов математического факультета. Какова в таком случае вероятность, что он окажется отличником? С точки зрения теории вероятностей, это будет \emph{условная вероятность события $A$ при условии $H_1$} (то есть $P(A|H_1)$). Нетрудно видеть, что эта вероятность в точности равна доле отличников среди студентов математического факультета (то есть 1/4).

Предположим, что мы знаем только те данные, которые приведены в условии задачи. Иными словами, нам даны вероятности событий $H_1$, $H_2$, $H_3$, а также условные вероятности события $A$ при условиях $H_1$, $H_2$, $H_3$ (то есть $P(A|H_1)$, $P(A|H_2)$, $P(A|H_3)$). Мы также знаем, что события $H_1$, $H_2$, $H_3$ попарно несовместны (то есть один студент может учиться только на одном факультете) и вместе образуют всё пространство элементарных исходов (любой студент учится хотя бы на одном факультете). Иными словами, эти события разбивают всё пространство элементарных исходов на несколько непересекающихся частей — в этом случае такие события часто называют \emph{гипотезами}, и говорят, что они образуют «полную систему событий». Тот факт, что гипотезы попарно несовместны, и что в объединении они дают всё пространство элементарных событий очень важен~--- он означает, что всегда реализуется в точности одна гипотеза.

Нетрудно видеть, что по указанным данным можно найти вероятность события $A$ (то есть долю отличников среди студентов всего вуза) — точно таким образом, как мы это сделали выше. Запишем приведенное решение в терминах теории вероятностей:

Напомним, что $P(A|H_1)$~--- это доля отличников на математическом факультете, $P(H_1)$~--- доля студентов математического факультета на всем вузе, и значит $P(A|H_1)\cdot P(H_1)$~--- доля отличников с математического факультета среди всех студентов вуза (см. формулу~\eqref{eq:doli}). Иными словами, это то же самое, что $P(A\cap H_1)$~--- вероятность одновременного выполнения $A$ и $H_1$ (в примере — вероятность выбрать одного из 25 отличников с математического факультета — она составляет 0.025). Аналогично $P(A|H_2)\cdot P(H_2)$~--- доля отличников со физического факультета среди всех студентов вуза, $P(A|H_3)\cdot P(H_3)$~--- доля отличников с химического факультета среди всех студентов вуза. Складывая эти доли, получаем долю всех отличников вуза, то есть $P(A)$.
\begin{equation}\label{eq:fp}
	P(A)=P(A|H_1)\cdot P(H_1)+P(A|H_2)\cdot P(H_2)+P(A|H_3)\cdot P(H_3),
\end{equation}

Говоря человеческим языком, это означает, что доля всех отличников в вузе это сумма долей отличников каждого факультета от всех студентов вуза.

\vspace{1em}
{\footnotesize
Это рассуждение можно записать более формально. Действительно, напомним, что по определению условной вероятности,
\begin{equation}
	P(A|H_1)=\frac{P(A\cap H_1)}{P(H_1)},
\end{equation}
где $P(A\cap H_1)$~--- вероятность одновременного выполнения $A$ и $H_1$. Умножая обе части этого равенства на $P(H_1)$, получим:
\begin{equation}\label{eq:AH1}
	P(A\cap H_1)=P(A|H_1)\cdot P(H_1).
\end{equation}
С другой стороны, поскольку события $H_1$, $H_2$, $H_3$ не пересекаются и вместе образуют всё пространство элементарных исходов, событие $A$ можно представить как объединение кусочков, лежащих в $H_1$, $H_2$, $H_3$, то есть $A=(A\cap H_1)\cup (A\cap H_2)\cup (A\cap H_3)$. События $A\cap H_1$, $A\cap H_2$, $A\cap H_3$ тоже не пересекаются, и значит их вероятности можно складывать: $P(A)=P((A\cap H_1)\cup (A\cap H_2)\cup (A\cap H_3))=P(A\cap H_1)+P(A\cap H_2)+P(A\cap H_3)$. Учитывая~(\ref{eq:AH1}), имеем: $P(A)=P(A|H_1)\cdot P(H_1)+P(A|H_2)\cdot P(H_2)+P(A|H_3)\cdot P(H_3)$. }
\vspace{1em}

Доказанная нами формула~(\ref{eq:fp}) называется \textit{формулой полной вероятности}. В общем случае гипотез может быть не три, а больше, и формула выглядит так.

Пусть имеется полная система событий $ {H_{1}, H_{2},..., H_{n}} $:~  все эти события попарно несовместны и их объединение составляет всё пространство элементарных исходов $\Omega$. Пусть нам известны вероятности $P(H_{1}),~P(H_{2}),\dots,~ P(H_{n})$~событий (гипотез) $H_{1},~ H_{2},\dots,~H_{n}$. Разумеется, $P(H_{1})+ P(H_{2})+\dots+ P(H_{n})=1$.  Пусть мы знаем условные вероятности выполнения события $A$ при условии выполнения каждого из событий $H_{1}, ~H_{2},\dots,~ H_{n}$, то есть $P(A|H_{1}),P(A|H_{2}),\dots, P(A|H_{n})$. Тогда вероятность события $A$ (полная вероятность) составляет

$$P(A)=P(H_{1})P(A|H_{1})+P(H_{2})P(A|H_{2})+\dots+P(H_{n})P(A|H_{n})$$



\textbf{Пример:}

1. На столе лежат три монетки: две обыкновенные и одна с орлами с двух сторон. Наудачу выбирается одна монетка и подбрасывается. С какой вероятностью выпадет орел?

Мы могли выбрать либо обычную монетку, либо монетку с двумя орлами. В каждом из этих случаев не сложно найти вероятность выпадения орла, поэтому обозначим за $H_{1}$ событие (гипотезу) ~--- мы взяли обычную монетку, за $H_{2}$ событие (гипотезу) ~--- мы взяли монетку с двумя орлами. Интересующее нас событие \textit{A} ~--- выпадение орла на выбранной монетке. Итак, имеем:\\

$P(H_{1})=2/3$,~~~~$P(A|H_{1})=1/2$

$P(H_{2})=1/3$,~~~~$P(A|H_{2})=1$  \\

Значит, $$P(A)=P(H_{1})P(A|H_{1})+P(H_{2})P(A|H_{2})=\frac{2}{3}\cdot\frac{1}{2}+\frac{1}{3}\cdot1=\frac{2}{3}.$$

\section{Формула Байеса}

Вернемся к задаче об отличниках. Пусть мы выбрали случайным образом студента, и он оказался отличником. Теперь нас интересует вероятность того, что он учится, скажем, на математическом факультете. Итак, задача в отыскании условной вероятности того, что студент учится на первом факультете, при условии того, что он отличник. Что мы знаем? Вероятности $P(H_{1})$, $P(H_{2})$, $P(H_{3})$  того, что студент учится на каждом из трех факультетов, вероятности (условные!) $P(A|H_{1}),~P(A|H_{2}),~P(A|H_{3})$ того, что студент, учась на каждом из этих факультетов, будет отличником. Нас же интересует условная вероятность $P(H_{1}|A)$, что студент, \textit{оказавшийся} отличником, учится на математическом факультете.

Иными словами, нам нужно вычислить, какую долю составляет площадь левого закрашенного прямоугольника на рис.~\ref{fig:bayes} (выделен более темным цветом), соответствующего отличникам с математического факультета (то есть событию $AH_1$), от площади всей закрашенной области (соответствующей всем отличникам, т.е.  событию $A$).

\begin{figure}[htb]
	\begin{center}
		\includegraphics[scale=0.7]{bayes.pdf}
		\caption{Формула Байеса: вычисление условной вероятности $H_1$ при условии $A$ соответствует вычислению доли площади левого тёмного прямоугольника во всей закрашенной области}\label{fig:bayes}
	\end{center}
\end{figure}

Напомним, что по формуле полной вероятности, $P(A)$ представляется в виде суммы трех слагаемых, каждое из которых выражает долю отличников с соответствующего факультета от числа всех студентов вуза (на рисунке — площади соответствующих закрашенных прямоугольников):
\begin{equation*}
	P(A)=P(A|H_1)\cdot P(H_1)+P(A|H_2)\cdot P(H_2)+P(A|H_3)\cdot P(H_3)
\end{equation*}
В этом месте есть кажущееся несоответствие условию задачи. Ведь мы \textit{знаем}, что студент оказался отличником, т.\,е.~ событие \textit{A} произошло, зачем же мы ищем вероятность этого события? На самом деле, здесь $P(A)$ \textit{априорная} вероятность события \textit{A}, просто вероятность выбрать отличника, выбирая случайного студента.

Нас же сейчас интересует, какую долю составляет первое слагаемое ($P(A|H_1)\cdot P(H_1)$) в этой сумме — это и есть условная вероятность выбрать отличника с математического факультета, выбирая из всех отличников. Иными словами, нас интересует отношение
\begin{equation*}
	P(H_1|A)=\frac{P(A|H_1)\cdot P(H_1)}{P(A)} = \frac{P(A|H_1)\cdot P(H_1)}{P(A|H_1)\cdot
	P(H_1)+P(A|H_2)\cdot P(H_2)+P(A|H_3)\cdot P(H_3)}
\end{equation*}
Вот мы и получили \emph{формулу Байеса} (в одной из своих записей). В рассмотренном примере интересующая нас вероятность есть отношение 25 отличников с математического факультета, ко всем 625 отличникам, то есть $25/625=0.04$ (или, переходя от студентов к долям от всего вуза, $0.025/0.625=0.04$).

\vspace{1em}
{\footnotesize Эти рассуждения также можно записать более формально. По определению условной вероятности	$$P(H_{1}|A)=\frac{P(H_{1}\cap A)}{P(A)}$$

Величину $P(A)$ мы нашли при помощи формулы полной вероятности 
$$P(A)=P(H_{1})P(A|H_{1})+P(H_{2})P(A|H_{2})+P(H_{3})P(A|H_{3}).$$

Осталось найти $P(H_{1}A)$. Это вероятность того, что случайно выбранный студент окажется отличником с математического факультета. Согласно~\eqref{eq:AH1}, $P(H_{1}A)=P(A\cap H_1)=P(A|H_{1})\cdot P(H_{1})$.

Итак, мы получили формулу Байеса:

\begin{equation*}
	P(H_1|A)=\frac{P(A|H_1)\cdot P(H_1)}{P(A)} = \frac{P(A|H_1)}{P(A|H_1)\cdot
	P(H_1)+P(A|H_2)\cdot P(H_2)+P(A|H_3)\cdot P(H_3)}
\end{equation*}

Она, разумеется, верна в случае любого количества гипотез:
	
\begin{equation}
	P(H_{i}|A)=\frac{P(H_{i})\cdot P(A|H_{i})}{P(H_{1})P(A|H_{1})+P(H_{2})P(A|H_{2})+...+P(H_{n})P(A|H_{n})}
\end{equation}
}
\textbf{Примеры:}

1. На столе лежат три монетки: две обыкновенные и одна с орлами с двух сторон. Наудачу выбирается одна монетка и подбрасывается. На подброшенной монетке выпал орел. С какой вероятностью мы кидали обычную монетку?

Пусть

$H_{1}$ событие (гипотеза) ~--- мы взяли обычную монетку,

$H_{2}$ событие (гипотеза) ~--- мы взяли монетку с двумя орлами,

Гипотезы $H_{1}$ и $H_{2}$ образуют полную систему событий.

Событие $A$~--- на монетке выпал орел.

Тогда
	
$P(H_{1})=2/3$, ~~~~
$P(A|H_{1})=1/2$

$P(H_{2})=1/3$, ~~~~
$P(A|H_{2})=1$

Значит,	$$P(A)=\frac{P(H_{1})P(A|H_{1})}{P(H_{1})P(A|H_{1})+P(H_{2})P(A|H_{2})}=\frac{\frac{2}{3}\cdot\frac{1}{2}}{\frac{2}{3}\cdot\frac{1}{2}+\frac{1}{3}\cdot1}=\frac{1/3}{2/3}=\frac{1}{2}$$

Это и неудивительно: хотя на странной монетке орлы встречаются в два раза чаще, чем на обычной, зато на нашем столе обычные монетки встречаются в два раза чаще, чем странные.

2. Тест на ВИЧ выдает верный результат в $95\%$ случаев. Известно, что ВИЧ заражена $1/1000$ часть всего населения. Человек получил положительный результат теста на ВИЧ (то есть по результатам теста он заражен). С какой вероятностью он действительно заражен?

Пусть

$H_{1}$ событие (гипотеза)  ~--- человек заражен ВИЧ,

$H_{2}$ событие (гипотеза) ~--- человек не заражен ВИЧ.

Гипотезы $H_{1}$ и $H_{2}$ образуют полную систему событий.

Событие $A$ ~--- тест дал положительный результат.

Тогда

$P(H_{1})=0,001$,~~~~      $P(A|H_{1})=0,95 $

$P(H_{2})=0,999$,~~~~      $P(A|H_{2})=0,05$\\

Значит,
$$	P(H_{1}|A)=\frac{P(H_{1})P(A|H_{1})}{P(H_{1})P(A|H_{1})+P(H_{2})P(A|H_{2})}=\frac{0,001\cdot0,95}{0,001\cdot0,95+0,999\cdot0,05}\approx0,0187$$

Результат кажется удивительным: ведь ошибка теста всего $5\%$, как же могло получиться, что вероятность заболевания при положительном результате теста так мала? Дело в том, что среди людей, получивших положительный результат теста, гораздо больше получивших ложный положительный результат ($5\%$ от $99,9\%$ населения), чем людей, действительно зараженных ВИЧ ($0{,}1\%$ населения), а тем более, действительно больных, получивших положительный результат теста~--- см. рисунок~\ref{fig:hivtest}.

\begin{figure}[htb]
	\begin{center}
		\includegraphics[scale=0.7]{hivtest.pdf}
		\caption{Применение формулы Байеса в случае теста на ВИЧ. Тонкая полоска слева~--- люди, зараженные ВИЧ. Темный закрашенный прямоугольник~--- те из них, кто получил положительный результат теста. Длинный закрашенный прямоугольник справа — здоровые люди, получившие положительный результат теста. Доля больных людей мала по сравнению со всеми людьми, получившими положительный результат.}\label{fig:hivtest}
	\end{center}
\end{figure}
\end{document}
